\documentclass[11pt,letterpaper]{article}

\usepackage[T1]{fontenc}
\usepackage{lmodern}
\usepackage[margin=1in]{geometry}
\usepackage{caption}
\usepackage[bookmarks=true]{hyperref}

% Reproducible output: no embedded build date/time or trailer id.
\pdfinfoomitdate=1
\pdftrailerid{}
\pdfsuppressptexinfo=-1

\pagestyle{empty}
\captionsetup{labelsep=colon}

\newcommand{\plot}{%
  \setlength{\unitlength}{1pt}%
  \begin{picture}(240,120)
    \put(0,0){\line(1,0){240}}
    \put(0,0){\line(0,1){120}}
    \put(0,0){\line(2,1){220}}
    \put(-6,-12){0}\put(110,-12){10}\put(225,-12){20}
    \put(120,60){latency}
  \end{picture}}

\begin{document}

\section{Introduction}

This document tests figure extraction and outlines. The text around each
figure is ordinary prose, long enough that every line holds many words.

\begin{figure}[h]
  \centering
  \plot
  \caption{Latency grows with load.}
\end{figure}

Figure 2 shows the same data from another angle, and this sentence names a
figure without being its caption, so it must not be read as one.

\subsection{Method details}

More prose follows here so that the section has a body of its own, with
several words on every line of the paragraph below the heading.

\begin{figure}[h]
  \centering
  \caption{Throughput by node, caption above.}
  \plot
\end{figure}

\section{Results}

The results section closes the document with one more paragraph of prose so
the outline has a second top-level entry.

\begin{thebibliography}{9}
\bibitem{smith} J. Smith and K. Lee, ``Fast storage for things,'' in
  \emph{Proc. FAST}, 2020, doi: 10.1145/12345.678.
\bibitem{doe} Jane Doe and John Roe. 2021. A Study of Caches. In
  \emph{Proceedings of SOSP}. https://example.org/caches
\bibitem{lee} Smith, J. and Lee, K. (2019). Learning to rank. arXiv preprint
  arXiv:1901.01234.
\end{thebibliography}

\end{document}
